% =====================================================================
% DRAFT v0.1 --- Multiview CRL on paired-modality brain MRI
% Status: initial skeleton. All quantitative results are from the
% synthetic pseudo-MRI benchmark. ADNI numbers are NOT yet in hand and
% are marked \pending throughout.
% =====================================================================
\documentclass[11pt]{article}

\usepackage[margin=1in]{geometry}
\usepackage{amsmath,amssymb,amsthm}
\usepackage{booktabs}
\usepackage{graphicx}
\usepackage{xcolor}
\usepackage{natbib}
\usepackage{hyperref}
\hypersetup{colorlinks=true,linkcolor=blue,citecolor=blue,urlcolor=blue}

% --- drafting macros -------------------------------------------------
\newcommand{\pending}[1]{{\color{red}\textbf{[PENDING: #1]}}}
\newcommand{\suspect}[1]{{\color{orange}#1}}
\newcommand{\todo}[1]{{\color{blue}\small[TODO: #1]}}
\newcommand{\zc}{\mathbf{z}_c}
\newcommand{\zs}{\mathbf{z}_s}
\newcommand{\R}{\mathbb{R}}

\title{\textbf{Contrastive Multiview Learning Re-Gauges Content Rather\\Than Identifying It: A Controlled Study on\\Paired-Modality Brain Imaging}}

\author{Natalia Glazman \\ \todo{affiliation, co-authors, supervisor}}
\date{Draft --- \today}

\begin{document}
\maketitle

\begin{abstract}
Multiview causal representation learning promises that content shared across
views is identifiable up to an invertible transformation, and paired-modality
medical imaging appears to be its ideal application: a T1- and a T2-weighted
scan of the same subject are two renderings of one underlying anatomy, differing
only in tissue contrast. We test this promise directly. Using a controlled
pseudo-MRI benchmark with a known causal generative process, we compare a
multiview contrastive objective against a pure reconstruction autoencoder
sharing an identical backbone, and evaluate both under a \emph{gauge-aware}
protocol that matches the probe's pooling to the identifiability claim being
made.

Our central finding is negative and, we argue, informative. The contrastive
objective does not recover more shared content than reconstruction: at the
pooling level where the two are fairly comparable, aggregate linear
recoverability is $0.409$ versus $0.409$, and peak recoverability across all
probe formats is tied to within $0.01$. The apparent contrastive advantage
reported by the field's default protocol --- a globally average-pooled linear
probe --- is largely a measurement artifact. We show by a no-retrain
intervention that the normalisation layer preceding the pooled probe discards
exactly the global channel statistics in which the reconstruction model stores
its content: the discarded statistics alone recover brain size at $R^2=0.68$,
better than the post-normalisation pooled representation they were removed from
($0.34$). Per-factor analysis converts this from artifact to mechanism: for
temporal atrophy, contrastive appears to win by $+0.31$ pooled, yet the
reconstruction baseline holds \emph{more} of that factor spatially ($0.682$
versus $0.623$). Contrastive learning relocates content from a spatial to a
global format --- it re-gauges --- and slightly degrades peak recoverability
while doing so.

What the contrastive objective does buy is style exclusion, and this is real but
partial: the dominant style factor still leaks into content at $0.481$, because
global intensity scale is entangled with global geometric scale, a content
factor the model must retain. It also carries a cost the pooled protocol cannot
see. Probing the latent grid voxel-wise, contrastive representations are
uniformly worse at localising content than the reconstruction baseline. For
medical imaging, where the quantity of interest is usually \emph{where} an
abnormality is and not merely whether a global embedding encodes it, this
trade-off is the operative one. We conclude with a gauge-aware evaluation
protocol, and argue that pooled linear probes systematically favour
invariance-based objectives over reconstruction-based ones for reasons unrelated
to identifiability.
\end{abstract}

% =====================================================================
\section{Introduction}
% =====================================================================

Representation learning for medical imaging is usually justified by an
identifiability argument, whether or not it is stated formally: we want the
learned latent variables to correspond to the physical or physiological factors
that generated the image, so that downstream analysis is measuring anatomy
rather than scanner. Recent multiview causal representation learning (CRL) makes
this argument precise. Given two views of the same underlying latent, generated
by distinct but invertible mixing functions, the \emph{shared} content block is
identifiable up to an invertible transformation of that block, under
contrastive-style objectives that align views while maintaining
entropy~\citep{yao2024multiview,vonkugelgen2021self,daunhawer2023identifiability}.

Paired-modality MRI looks like the canonical instance of this setting. A
T1-weighted and a T2-weighted scan of the same session share anatomy exactly ---
they are, after registration, two functions of one tissue geometry --- and
differ in a low-dimensional contrast transformation determined by pulse sequence
and scanner. Content is anatomy; style is contrast. Unlike the augmentation-based
views of self-supervised learning, where the ``shared content'' is defined
implicitly and somewhat circularly by the augmentation policy, here the split is
physically grounded.

This paper reports what happened when we took that correspondence seriously and
measured it. Our conclusion is that the identifiability framing, applied
naively, mismeasures what these objectives do. Three findings drive this:

\begin{enumerate}
\item \textbf{Content completeness is tied.} A multiview contrastive objective and
a pure reconstruction autoencoder with the same backbone recover the same amount
of ground-truth shared content, when both are read out at a pooling level that
does not privilege one storage format (Section~\ref{sec:completeness}).

\item \textbf{The reported advantage is a gauge artifact.} The default protocol ---
globally average-pooled linear probe --- reports a contrastive win. We show by
intervention that this win is produced by a normalisation layer that deletes the
global channel statistics in which the reconstruction model stores content
(Section~\ref{sec:gauge}). Reconstruction is not worse; it is being read in the
wrong basis.

\item \textbf{The genuine difference is a localisation--isolation trade-off.}
Contrastive learning does exclude style from content more effectively, but it
delocalises: at the voxel level it is uniformly worse at saying \emph{where} a
factor is (Section~\ref{sec:voxel}). For clinical use this is the trade-off that
matters, and the pooled protocol is blind to it.
\end{enumerate}

We frame these through a distinction we make explicit in
Section~\ref{sec:gauge-theory}: between \emph{block identifiability}, which
holds content up to an arbitrary invertible map and is therefore agnostic to
spatial format, and what we call \emph{voxel-compositional identifiability},
which additionally requires content to be recoverable at the spatial location it
describes. Multiview CRL theory guarantees the first. Medical imaging
applications almost always want the second. The two come apart precisely here,
and the metrics used to certify the first are systematically miscalibrated for
comparing methods that differ in the second.

\paragraph{Contributions.}
\begin{itemize}
\item A controlled pseudo-MRI benchmark with a known causal content SCM,
view-specific style, and --- importantly --- an explicit account of its own
unlabeled shared nuisance dimensions (Section~\ref{sec:benchmark}).
\item A gauge-aware evaluation protocol: pooling-matched probes, null-subtracted
leakage, separated completeness and isolation axes, and voxel-level factor
localisation (Section~\ref{sec:protocol}).
\item The empirical characterisation above, with a mechanistic account of each
result rather than a leaderboard (Sections~\ref{sec:results}--\ref{sec:why}).
\item Practical guidance for evaluating disentanglement in paired-modality
imaging, including two failure modes of patch-level contrastive objectives on
spatially registered data (Section~\ref{sec:why}).
\end{itemize}

% =====================================================================
\section{Background: identifiability and its gauge}
\label{sec:gauge-theory}
% =====================================================================

\paragraph{Setup.} Let $\zc \in \R^{d_c}$ be content and $\zs^{(k)} \in
\R^{d_s}$ be view-specific style for view $k \in \{1,2\}$. Observations are
$x^{(k)} = g_k(\zc, \zs^{(k)})$ for injective mixing $g_k$. An encoder
$f_\theta$ produces $\hat{\zc}$; the multiview result states that under
suitable conditions the contrastive optimum satisfies $\hat{\zc} = h(\zc)$ for
some invertible $h$, with $\hat{\zc}$ independent of $\zs^{(k)}$.

\paragraph{The gauge group.} The theorem's guarantee is an equivalence class,
not a representation. Any $h \in \mathrm{GL}$-like invertible family yields an
equally valid solution. We call this the \emph{gauge freedom} of the claim, and
two solutions related by such an $h$ are \emph{gauge-equivalent}. This matters
for evaluation in a way that is easy to overlook:

\begin{quote}
\emph{A metric is only meaningful for an identifiability claim if it is
invariant to that claim's gauge group.}
\end{quote}

A linear probe on globally average-pooled features is \emph{not} invariant to
the gauge group of block identifiability. Two gauge-equivalent encoders --- both
perfectly block-identified --- can score arbitrarily differently under it if
$h$ redistributes information between spatial layout and channel statistics.
This is not a hypothetical: it is exactly what we measure in
Section~\ref{sec:gauge}.

\paragraph{Two identifiability notions.} We distinguish:

\begin{description}
\item[Block identifiability (BI).] $\hat{\zc} = h(\zc)$, $h$ invertible. Gauge
group: all invertible maps on the content block. Spatially blind. This is what
the multiview theory delivers.

\item[Voxel-compositional identifiability (VCI).] The representation is a field
$\hat{c}(v)$ over spatial locations $v$, and content factors are recoverable
from the neighbourhood of the locations they physically affect. Gauge group:
permutation/invertible maps \emph{within} location, not across. Strictly
stronger than BI.
\end{description}

VCI is the informal target of most medical imaging work --- lesion
localisation, voxel-based morphometry, regional atrophy --- and it is what
object-centric identifiability results in vision aim
at~\citep{brendel2023objectcentric}. To our knowledge no identifiability theory
covers the content/style structure of paired-modality anatomical imaging, and
establishing that theory is beyond this paper's scope. Our contribution is to
show empirically that BI and VCI dissociate here, in a direction that matters,
and that the standard protocol reports only BI while practitioners read it as
VCI. \todo{Decide how hard to push the VCI formalism --- either state it as a
definition with a proposition about the gauge group, or keep it informal and
purely as an evaluation-design argument. Currently informal.}

% =====================================================================
\section{A controlled pseudo-MRI benchmark}
\label{sec:benchmark}
% =====================================================================

Real paired MRI has no ground-truth latents, so all identifiability claims on it
are indirect. We therefore built a synthetic generator
(\texttt{eval/synthetic\_dataset.py}) that renders two contrast-differing
volumes from one explicit anatomical latent.

\paragraph{Content.} Nine named scalar factors drive a 3D tissue-label
geometry: brain size, ventricle size, cortical thickness, temporal-lobe
atrophy, left--right asymmetry, sulcal widening, and a three-dimensional
white-matter lesion position. Factors may be sampled i.i.d.\ or from a linear
non-Gaussian structural causal model over a chain/random DAG
(\texttt{build\_content\_scm}), so that content factors are \emph{dependent} ---
the regime the multiview theory is designed for and where component-wise
disentanglement metrics are inappropriate.

\paragraph{Style.} Three view-specific scalars --- intensity gain, bias, and
noise $\sigma$ --- are drawn \emph{independently per view}
($\zs^{(1)} \perp \zs^{(2)}$), then applied together with a modality-specific
tissue-intensity lookup table (T1: WM $0.8$, GM $0.5$; FLAIR: WM $0.4$,
GM $0.8$) and a smooth multiplicative bias field, followed by Rician noise. The
independent-per-view draw is deliberate: it means the contrastive objective has
full leverage to remove style, since style genuinely differs between the aligned
positive pair. Any failure to exclude style is therefore a property of the
objective, not a structural impossibility.

\paragraph{The benchmark's own confound (and why we report it).} The renderer
also perturbs tissue boundaries with two unlabeled, view-invariant random
fields: a deformation field on a $4^3$ grid ($64$ dimensions) and a fissure
field on an $8^3$ grid ($512$ dimensions). These are \emph{shared content} by
every definition --- identical across views, subject-specific --- but they are
not in the probed ground truth. True shared content is therefore
$\approx\!585$-dimensional while the evaluation probes $9$ of it.

This has two consequences we believe generalise beyond our generator, and which
we report rather than quietly fix:

\begin{enumerate}
\item \textbf{Aggregate factor-recovery ceilings are not what they appear.} An
aggregate $R^2 \approx 0.35$ across the nine factors reads as a model failure
but is largely a coverage artifact: the probe measures under $2\%$ of the shared
content the model correctly encodes.
\item \textbf{It creates a contrastive shortcut.} A $585$-dimensional
view-invariant, subject-unique signature solves instance discrimination on its
own. At batch size $B$, InfoNCE needs only $\log_2 B$ bits ($6$ bits at
$B{=}64$), and the corrugation field supplies them. Top-1 accuracy saturates
early and no further content is driven into the representation. Critically,
\emph{scaling the task does not fix this} --- larger batches and harder
negatives simply discriminate the same signature more sharply. Because these
fields are trilinear upsamples of coarse grids, they are smooth and
volume-spanning by construction, so an encoder that solves InfoNCE through them
necessarily produces a low-frequency, spatially diffuse content map. Some of the
delocalisation we report in Section~\ref{sec:voxel} is the model faithfully
encoding a delocalised signal.
\end{enumerate}

We provide a \texttt{clean\_content} mode that zeroes both fields, isolating the
nine labeled factors. \todo{Report the main comparison under BOTH modes. Current
headline numbers are from the nuisance-on setting; the clean-mode replication is
the single most valuable missing experiment.}

\paragraph{Resolution caveat.} The lesion is a radius-$0.1$ sphere in
$[-1,1]^3$ at $64^3$, i.e.\ $\approx\!3.2$ voxels in radius. At the $16^3$
latent grid it occupies $\approx\!1.6$ latent voxels: resolvable in principle,
but averaged away by pooled probes ($\approx\!3\%$ of a $4{\times}4{\times}4$
patch, $0.4\%$ of a $2{\times}2{\times}2$ eval patch). Near-zero lesion recovery
in pooled readouts is therefore \emph{not} evidence of absence
(Section~\ref{sec:limitations}).

% =====================================================================
\section{Models and objectives}
\label{sec:models}
% =====================================================================

All arms share one backbone so that the objective is the only manipulated
variable: a hierarchical 3D VQ-VAE-2~\citep{razavi2019vqvae2} with three levels.
At the finest level, hidden channels are partitioned into content and style
groups by a mask (learned via Gumbel-softmax straight-through, or fixed
first-$K$); content channels feed the codebook and the contrastive loss, style
channels feed decoder injection (concatenation or spatial FiLM) and an optional
independent style codebook. A \texttt{SplitGroupNorm} normalises content and
style groups independently so that modality-specific style statistics cannot
bias content activations before the mask --- a design choice that turns out to
be central to Section~\ref{sec:gauge}.

\paragraph{Arms.}
\begin{description}
\item[\textsc{Recon}] Reconstruction only. No contrastive term. The autoencoder
baseline that multiview CRL is claimed to beat.
\item[\textsc{Contrast}] Reconstruction plus a multiview contrastive term on
content channels (InfoNCE by default; Barlow Twins and VICReg also implemented),
at a $10{:}1$ contrastive-to-reconstruction weighting. Cross-view-negatives-only
restricts the InfoNCE denominator to the opposite view.
\item[\textsc{EncOnly}] \pending{not yet run} An encoder-only arm with no
decoder, matching the multiview CRL formulation exactly (alignment minus
entropy). Implemented in \texttt{models/multiview\_encoder.py} +
\texttt{training/main\_conv\_synthetic.py}, reusing the same encoder stack for a
controlled reconstruction-vs-no-reconstruction ablation.
\end{description}

\paragraph{A fidelity note.} The multiview CRL method is encoder-only:
invertibility is enforced by an entropy term, not a decoder, and reconstruction
is the \emph{baseline} it outperforms~\citep{yao2024multiview}. Our
\textsc{Contrast} arm therefore is not a reproduction of that method but a test
of whether its contrastive term adds identifiability \emph{on top of} a
reconstruction model --- which is the realistic setting for medical imaging,
where a decoder is usually wanted anyway. The \textsc{EncOnly} arm closes this
gap and is the highest-priority missing experiment.

\paragraph{No projection head.} Following the multiview CRL formulation, the
contrastive loss is applied directly to content channels with no projection MLP,
so the loss shapes exactly the representation that is subsequently probed. This
departs from the SimCLR/MoCo/VICReg convention of interposing a
2--3 layer head and probing pre-head features, a convention motivated precisely
by the observation that the loss-facing space over-compresses toward invariance
and loses linearly decodable information. We flag this as a candidate partial
explanation for the tie in Section~\ref{sec:completeness}; a projection-head
ablation is implemented (\texttt{--contrastive-proj-dim}) but
\pending{not yet run}.

% =====================================================================
\section{Gauge-aware evaluation protocol}
\label{sec:protocol}
% =====================================================================

Our protocol separates two axes that the literature routinely reports as one
number, and matches probe pooling to the claim.

\paragraph{Two axes.}
\begin{description}
\item[Completeness.] How much of $\zc$ is recoverable from $\hat{\zc}$. Measured
by cross-validated probe $R^2$ (RidgeCV linear and a two-layer MLP, $5$-fold,
multi-seed) and by \emph{block-MCC}: fit a readout $\hat{\zc} \to \zc$,
Hungarian-match predicted to true dimensions on $|\rho|$, average the matched
correlations. Block-MCC is the correct MCC under BI's gauge group; per-channel
MCC is not, and we do not report it as a headline.
\item[Isolation.] How much of $\zs$ leaks into $\hat{\zc}$. Measured by
null-subtracted content$\to$style probe $R^2$, and by
$\mathrm{separation} = \mathrm{MCC}(\hat{c}\to\zc) - \mathrm{MCC}(\hat{c}\to\zs)$.
\end{description}

\paragraph{Pooling-matched readout.} Every probe is run at three poolings ---
global average (GAP), per-channel statistics, and a spatial patch grid --- and
we report the \emph{best over poolings} per factor per model as the
gauge-appropriate summary of BI completeness. Reporting a single pooling
silently fixes a gauge and, as we show, decides the result.

\paragraph{Null subtraction.} Every metric is accompanied by a
label-permutation null computed under identical probe, pooling, and split, and
we report the null-subtracted gap. This matters most for high-dimensional patch
features where flat probes overfit.

\paragraph{Voxel-level localisation.} For VCI, we probe the native latent grid
per location: a shared probe predicts a ground-truth field (tissue class, lesion
mask, deformation) from each latent voxel's feature vector, against a
subject-shuffled null. We summarise with peak in-brain AUC/$R^2$, and with
\emph{in-mass}: the fraction of above-null predictive mass falling inside the
region the factor physically affects. Uniform mass would be $\approx\!0.145$.

\paragraph{Controls.} Three controls gate every conclusion: (i) a positive
control (tissue class must be decodable, verifying ground-truth$\leftrightarrow$
feature alignment), (ii) a raw-pixel readout floor, and (iii) an effective-rank
measure (Roy--Vetterli) to rule out dimensional collapse as an explanation for
probe differences.

% =====================================================================
\section{Results}
\label{sec:results}
% =====================================================================

Unless stated otherwise, all results are on the synthetic pseudo-MRI benchmark
with causal (dependent) content, comparing \textsc{Contrast} against
\textsc{Recon} at matched backbone and data. \pending{ADNI replication --- the
training and evaluation pipeline is complete and the separation-score /
recovery-map machinery is validated on synthetic, but no ADNI identifiability
numbers are yet in hand. Every claim below is synthetic-only.}

% ---------------------------------------------------------------
\subsection{The default protocol reports a contrastive win}
% ---------------------------------------------------------------

Read out the way the field usually reads out --- linear probe on
globally-average-pooled content --- and contrastive learning wins. Aggregate
linear $R^2$ improves by $+0.060$ at GAP. The gain is concentrated in subtle,
spatially distributed factors: temporal atrophy $+0.31$ (about $64\%$ of the
aggregate gap), sulcal widening $+0.12$, left--right asymmetry $+0.10$,
ventricle size $+0.07$, with ties on saturated factors (brain size, cortical
thickness, both $\approx\!0.79$).

Taken alone this is a clean positive result, and it is the result we would have
reported under the standard protocol. The rest of this section is why it does
not survive.

% ---------------------------------------------------------------
\subsection{The win vanishes under pooling-matched readout}
\label{sec:completeness}
% ---------------------------------------------------------------

The advantage shrinks monotonically as pooling becomes more spatial: $+0.060$ at
GAP, $+0.015$ at channel statistics, $+0.010$ at patch. At the patch level, peak
recoverability is tied to three decimal places (Table~\ref{tab:completeness}).

\begin{table}[t]
\centering
\caption{Content completeness, \textsc{Contrast} vs \textsc{Recon}. Aggregate
over the nine labeled content factors. The contrastive advantage exists only at
GAP.}
\label{tab:completeness}
\begin{tabular}{llccc}
\toprule
Readout & Pooling & \textsc{Contrast} & \textsc{Recon} & $\Delta$ \\
\midrule
Linear $R^2$        & GAP    & --- & --- & $+0.060$ \\
Linear $R^2$        & stats  & --- & --- & $+0.015$ \\
Linear $R^2$        & patch  & $0.398$ & $0.388$ & $+0.010$ \\
Linear $R^2$ (mean, per-factor) & patch & $0.409$ & $0.409$ & $0.000$ \\
MLP $R^2$           & patch  & $0.289$ & $0.334$ & $-0.045$ \\
MLP $R^2$ (best pooling) & ---   & $0.323$ (gap) & $0.334$ (patch) & $-0.011$ \\
Block-MCC           & patch  & $0.540$ & $0.535$ & $+0.005$ \\
Best over poolings (mean) & --- & $0.419$ & $0.416$ & $+0.003$ \\
\bottomrule
\end{tabular}
\todo{Fill the two missing absolute GAP/stats cells from
\texttt{content\_rank\_out6}; add seed-level error bars (seeds 0--4 exist).}
\end{table}

Two secondary observations rule out the obvious alternative explanations. The
MLP probe never exceeds the linear probe in any cell ($-0.02$ to $-0.11$), so
there is no nonlinearly-hidden content that the linear probe misses --- and
since the MLP uses early stopping and thus trains on less data, the honest
reading is $\text{MLP} \approx \text{linear}$, i.e.\ zero nonlinear advantage.
And effective rank rules out dimensional collapse: \textsc{Contrast} has
\emph{higher} GAP effective rank ($29$ vs $13$) and patch rank is comparable
($115$ vs $127$). The reconstruction model's low GAP rank is spatial storage,
not collapse.

A raw-pixel control confirms both models are doing real work: a $16^3$
raw-intensity readout is strongly negative (linear $-0.44$, MLP $-2.5$) against
model $R^2 \approx 0.36$. The encoders extract content that pixels do not
linearly expose, as expected given inverted T1/FLAIR contrast.

% ---------------------------------------------------------------
\subsection{The GAP win is a normalisation artifact}
\label{sec:gauge}
% ---------------------------------------------------------------

Why is the reconstruction model so much worse at GAP specifically? A no-retrain
intervention answers this. We forward-hook \texttt{SplitGroupNorm} and probe
content features immediately before and after it
(\texttt{eval/probe\_prenorm\_groupnorm.py}):

\begin{itemize}
\item GAP recovery roughly \textbf{halves} across the layer: mean $R^2$
$0.34 \to 0.18$; brain size $0.78 \to 0.34$; left--right asymmetry
$0.59 \to 0.21$.
\item Patch recovery is \textbf{unaffected}: $0.307 \to 0.309$. The norm
preserves spatial pattern and destroys only the global per-channel statistic.
\item \textbf{Smoking gun:} the discarded per-group $(\mu,\sigma)$ statistics
\emph{alone} recover brain size at $R^2 = 0.68$ --- nearly double the $0.34$
obtainable from the post-normalisation pooled features. The factor literally
lives in what the normalisation subtracts.
\end{itemize}

This is a gauge effect in the precise sense of
Section~\ref{sec:gauge-theory}. The reconstruction model stores global content
in per-channel scale; normalisation deletes that channel; GAP pooling has
nothing left to read. The contrastive model, trained on normalisation-immune
patch features, stores the same content in a format that survives. Both encode
it. Only one is legible to the standard probe.

The practical consequence is a claim about the field's methodology, not just
about our models: \textbf{globally-pooled linear probes are not a neutral
instrument for comparing invariance-based against reconstruction-based
objectives}, because the two systematically differ in the representational
format the probe is sensitive to. \todo{Check how widely SplitGroupNorm-like
normalisation appears in comparable SSL pipelines --- this determines whether
the methodological claim generalises or is specific to our backbone. Important
for the framing; currently overstated relative to the evidence.}

% ---------------------------------------------------------------
\subsection{Re-gauging, proven per factor}
% ---------------------------------------------------------------

Aggregates could hide a genuine mixed result, so we analyse per factor,
comparing each model's \emph{best} recoverability across all four
(probe $\times$ pooling) formats rather than its GAP-linear score
(seeds $0$--$4$).

The decisive case is temporal atrophy --- the factor responsible for most of the
headline GAP gain:

\begin{center}
\begin{tabular}{lcc}
\toprule
Temporal atrophy, $R^2$ & \textsc{Contrast} & \textsc{Recon} \\
\midrule
Linear, GAP   & $\mathbf{0.486}$ & $0.173$ \\
Linear, patch & $0.623$ & $\mathbf{0.682}$ \\
\bottomrule
\end{tabular}
\end{center}

The direction \emph{reverses} with pooling. Contrastive learning does not add
temporal-atrophy information; the reconstruction model holds more of it, stored
spatially, where the normalisation-plus-GAP path cannot see it. The apparent
$+0.31$ is a patch$\to$GAP relocation that slightly \emph{destroys} peak
content.

The pattern holds across factors. On best-format recovery, $\Delta$
(\textsc{Contrast} $-$ \textsc{Recon}) is negative or flat for most factors
(temporal atrophy $-0.06$, cortical thickness $-0.04$, left--right asymmetry
$-0.02$, brain size $-0.02$); only sulcal widening ($+0.13$) and ventricle size
($+0.06$) are genuine peak gains; lesions are $\approx\!0$ in both. Aggregate
best-format means are $0.419$ vs $0.416$.

\paragraph{Summary claim.} Contrastive learning does not improve content
identifiability in the BI sense --- peak readout is tied --- but it
\emph{reformats} subtle anatomical content from a spatially distributed encoding
into a globally linearly-decodable one, at a small cost to peak recoverability
for the most spatially-stored factors. In the language of
Section~\ref{sec:gauge-theory}: it selects a different representative within the
same equivalence class.

% ---------------------------------------------------------------
\subsection{Isolation: a real gain, but partial}
\label{sec:isolation}
% ---------------------------------------------------------------

Completeness is only half of block identifiability; the other half is isolation
--- $\hat{\zc}$ free of style --- and it is the half InfoNCE's mechanism
actually targets. Here contrastive learning does win, but not cleanly.

Qualitatively the win is clear: the reconstruction baseline shows symmetric
style leakage into content ($+0.16$/$+0.28$ on the leakage diagonal) while the
contrastive model's is clean, and content$\to$style leakage is much lower
($0.18$ vs $0.49$). \suspect{However, the headline separation comparison
($0.158$ vs $-0.087$) is confounded: the baseline was scored with all channels
treated as content, so its ``leakage'' counts style channels as content. The
per-block re-run is required before any of these numbers are reportable.}
\pending{fair \texttt{--baseline-per-block} re-run --- highest-priority fix}

The contrastive model's \emph{own} per-latent numbers are valid, and they
complicate the story. Style has three factors, but only gain is substantively
encoded at all (style$\to$style: gain $0.522$, bias $0.03$, noise $0.01$ --- the
latter two are too weak to be represented). And gain leaks into content at
$\mathbf{0.481}$. The low aggregate leak of $0.18$ is therefore misleading: it
averages a genuine failure on the one real style factor against two factors
neither model encodes. Symmetrically, the style block encodes brain size at
$0.60$.

The mechanism is a same-character entanglement. Gain is a global \emph{intensity}
scale; brain size is a global \emph{geometric} scale. The encoder must retain
brain size (it is content), and it apparently cannot separate the two global
scale signals, so a shared global-scale feature carries both. We emphasise that
this is not a structural impossibility: gain is drawn independently per view, so
it differs across the aligned positive pair and the contrastive objective has
full leverage to remove it. This is a genuine isolation failure.

\paragraph{Writeable finding.} \emph{Even with view-varying style and an explicit
contrastive objective, the dominant style factor is not isolated, because it is
entangled with a content factor of the same character.} We suggest this
generalises: content/style splits are hard exactly where the two blocks contain
same-type quantities, and paired-modality MRI has this property (global
intensity scaling versus global head size; regional contrast versus regional
tissue change).

% ---------------------------------------------------------------
\subsection{The voxel-level trade-off}
\label{sec:voxel}
% ---------------------------------------------------------------

Pooled metrics cannot see \emph{where} content lives. Probing the native latent
grid per location tells a consistent and, for medical imaging, decisive story:
\textbf{contrastive learning uniformly degrades localisation}.

\begin{table}[t]
\centering
\caption{Voxel-level factor localisation on the native $16^3$ latent grid,
shared probe vs subject-shuffled null, $300$ samples.}
\label{tab:voxel}
\begin{tabular}{lccc}
\toprule
Target & \textsc{Contrast} & \textsc{Recon} & Null \\
\midrule
Tissue (GM), AUC \emph{(positive control)} & $0.74$ & $\mathbf{0.88}$ & --- \\
Lesion, AUC & \suspect{$0.78$} & \suspect{$\mathbf{0.95}$} & $0.46$/$0.57$ \\
Deformation, $R^2$ & $0.007$ & $\mathbf{0.061}$ & --- \\
\midrule
Cross-view content consistency & $\mathbf{0.65}$ & $0.47$ & --- \\
Style leakage & \textbf{clean} & $+0.16$/$+0.28$ & --- \\
\bottomrule
\end{tabular}
\end{table}

Every localised target is worse under the contrastive objective, while
cross-view consistency and style purity are better. This is an explicit
trade-off, not a capacity limit: the lesion is voxel-identifiable in
\emph{both} models (well above null in each), so it is not being quantised away.
The global contrastive objective buys view-invariance and style exclusion by
sacrificing local content fidelity --- it discards local content along with the
style.

A delocalisation signature supports this reading: content migrates outward, with
brain size recovered better \emph{outside} the brain ($0.80$) than inside
($0.60$), i.e.\ into background registers.

\paragraph{Two honest qualifications.} First, \suspect{the lesion AUC numbers are
suspect}: the evaluated checkpoint predates a generator fix that changed lesion
rendering, so it was scored on out-of-distribution images. Tissue, brain size,
and style findings are unaffected (those factors were unchanged), but the
headline lesion gap needs re-running.
\pending{lesion re-run under matched generator vintage}

Second, on a generator-matched re-run of the broad morphometric factors
(temporal atrophy, asymmetry, sulcal widening), the clean ``baseline localises,
contrastive smears'' pattern only \emph{partially} replicates. Both models are
heavily outside-biased (in-mass $0.02$--$0.18$ against a uniform $0.145$), with
contrastive roughly twice as outside-biased ($0.07$ vs $0.16$--$0.18$) --- the
same direction, but modest. The reason is instructive: those factors were, in
that generator vintage, spatially broad (half-brain, whole-hemisphere), and a
broad factor is legitimately decodable from boundary voxels' receptive fields.
\textbf{Voxel localisation is only a well-posed question for focal factors.}
Demonstrating it cleanly requires a focal factor and a generator-matched model,
which is our primary planned experiment (Section~\ref{sec:next}).

% =====================================================================
\section{Why: mechanisms}
\label{sec:why}
% =====================================================================

We collect the mechanistic account, since we think it is more transferable than
the numbers.

\paragraph{1. The objective has no locality term.} Global-pooled contrastive
learning rewards the \emph{pooled} content for being view-invariant and
discriminable. Nothing in the objective prefers that content sit at the location
it describes. Delocalisation is therefore not a bug but the optimum of a
localisation-blind objective.

\paragraph{2. Invariance is not informativeness.} InfoNCE optimises instance
discrimination. Once the task is solved --- and Section~\ref{sec:benchmark}
shows it is solved early and cheaply via a high-dimensional shared nuisance
signature --- gradients vanish and subtle factors are never driven into content.
Reconstruction, by contrast, must encode every voxel-moving factor. Hence
$\textsc{Recon} \geq \textsc{Contrast}$ on completeness is the \emph{predicted}
ordering, not an anomaly.

\paragraph{3. Patch-level InfoNCE has structural false negatives on registered
data.} Spatial registration makes same-location patches from different subjects
coincide by construction. At uniform interior or background patches, the only
signal separating ``negatives'' is per-subject nuisance or style --- so InfoNCE
is forced to inject nuisance and style \emph{into} content channels. We measured
this: on a $4^3$ grid, $42\%$ of patches are dead background at the batch size
actually used ($B{=}32$; $23\%$ at $B{=}128$, $16\%$ at $B{=}400$), and of the
remainder roughly a third are uniform-interior false-negative zones. Global
pooling launders this problem --- whole volumes differ in aggregate content, so
negatives are clean --- which is why GAP contrastive reaches clean linear
separation while patch block-MCC stays tied to the baseline.

The prescription follows: use negatives where they are clean (InfoNCE at GAP)
and negative-free objectives where they are not (Barlow Twins / VICReg on
\emph{foreground} patches). A negative-free loss is trivially aligned at a
uniform patch --- no gradient, no harm --- whereas InfoNCE must discriminate and
therefore does harm. Patch-level readout remains perfectly sound as a
\emph{probe}; it is only optimisation with negatives that is unsafe.

\paragraph{4. No projection head.} Applying the loss directly to probed features
means the invariance-induced over-compression is measured rather than shielded
(Section~\ref{sec:models}). This is faithful to the multiview CRL formulation
but confounds the comparison with the standard SSL recipe, and is a live
alternative explanation for part of the completeness tie.

% =====================================================================
\section{Limitations and threats to validity}
\label{sec:limitations}
% =====================================================================

We state these plainly; several materially qualify the results.

\begin{enumerate}
\item \textbf{No real-data results.} Every quantitative claim is on the
synthetic benchmark. The ADNI pipeline is complete and the evaluation machinery
validated, but no ADNI identifiability numbers exist yet. The paper's scope is
currently ``a controlled study'', not ``a study of brain MRI''.
\pending{ADNI replication of the completeness tie and the localisation gap}

\item \textbf{The isolation headline is confounded.} The baseline was scored
without a content/style split, inflating its apparent leakage. The
separation numbers in Section~\ref{sec:isolation} must not be reported until the
per-block re-run is done.

\item \textbf{The lesion localisation gap is on out-of-distribution images} for
the evaluated checkpoint (generator-fix boundary). Directionally consistent with
the other factors, but not yet trustworthy on its own.

\item \textbf{Near-zero lesion recovery is partly a probe artifact.} At a finer
$8^3$ eval grid the flat probe overfits badly (feature dimension $22{,}528$
against $\approx\!2{,}000$ samples), driving lesion $R^2$ \emph{more} negative;
PCA truncation confirms this (patch $R^2$ at $24$ components $0.325$ exceeds
full-block $0.228$). The lesion encodability question is genuinely open and
needs a dimensionality-controlled probe.

\item \textbf{Single split on voxel probes.} The voxel-level results use one
$70/30$ split without seed replication. Error bars needed before publication.

\item \textbf{Single backbone, single dataset, largely single seed} on the main
comparison. The completeness tie is replicated across seeds $0$--$4$ for
per-factor analysis, but the architectural generality of the gauge argument is
untested.

\item \textbf{The encoder-only arm is missing}, so we have not tested the actual
multiview CRL method --- only the value of adding its contrastive term to a
reconstruction model.
\end{enumerate}

% =====================================================================
\section{Planned experiments}
\label{sec:next}
% =====================================================================

The decisive experiment is an ablation grid trained on the current (focal-factor)
generator: $\{$\textsc{Recon}, global InfoNCE, foreground-masked patch
BT/VICReg, $+$ detached content injection$\}$, each evaluated with voxel-level
localisation, null-subtracted style leakage, and recovery maps.

\paragraph{Pre-registered outcome.} If foreground-masked patch objectives raise
focal-factor in-brain localisation toward the reconstruction baseline
\emph{while} keeping style leakage at or below zero, the trade-off is breakable
and we have a positive methodological result. If localisation stays flat, the
trade-off is fundamental to invariance-based objectives on registered imaging,
and we have a characterisation result. Both are publishable; we commit to the
criterion in advance to avoid selecting the framing post hoc.

Secondary: the \texttt{clean\_content} replication (removing the nuisance
shortcut), the encoder-only arm, the projection-head ablation, and the fair
per-block isolation scoring.

% =====================================================================
\section{Conclusion}
% =====================================================================

Multiview contrastive learning is often adopted in medical imaging on the
strength of an identifiability guarantee. We find that on paired-modality data
the guarantee is being measured with an instrument that does not respect its own
gauge group, and that the resulting comparisons systematically favour
invariance-based objectives over reconstruction-based ones for reasons unrelated
to how much content is identified.

When measured in a gauge-appropriate way, a multiview contrastive term added to
a reconstruction autoencoder does not identify more shared content. What it does
is change the format in which content is stored --- from spatially distributed
to globally linearly-decodable --- and exclude style from content more
effectively, though incompletely, failing precisely where a style factor shares
its character with a content factor. It pays for this with spatial localisation.

For applications that want a single style-invariant global embedding per
subject, that is a good trade. For applications that want to know where the
pathology is --- most of medical imaging --- it is currently a bad one, and a
plain reconstruction model localises better. Whether the two goals can be
satisfied simultaneously, by applying invariance locally rather than globally,
is the question our next experiment is designed to answer.

% =====================================================================
\section*{Reproducibility}
% =====================================================================
Code, synthetic generator, and evaluation protocol:
\todo{repository URL}. All evaluation entry points
(\texttt{eval/run\_dci\_compare.py}, \texttt{eval/content\_rank\_pca.py},
\texttt{eval/probe\_prenorm\_groupnorm.py}, \texttt{eval/recovery\_maps.py})
report null-subtracted metrics with fixed seeds. Because generator semantics
changed during the project, every reported result should be paired with the
generator commit SHA used at training time.

\bibliographystyle{plainnat}
\begin{thebibliography}{99}

\bibitem[Yao et al.(2024)]{yao2024multiview}
D.~Yao, D.~Xu, S.~Lachapelle, S.~Magliacane, P.~Taslakian, G.~Martius,
J.~von K\"ugelgen, and F.~Locatello.
\newblock Multi-view causal representation learning with partial observability.
\newblock \emph{ICLR}, 2024. arXiv:2311.04056.

\bibitem[von K\"ugelgen et al.(2021)]{vonkugelgen2021self}
J.~von K\"ugelgen, Y.~Sharma, L.~Gresele, W.~Brendel, B.~Sch\"olkopf,
M.~Besserve, and F.~Locatello.
\newblock Self-supervised learning with data augmentations provably isolates
content from style.
\newblock \emph{NeurIPS}, 2021.

\bibitem[Daunhawer et al.(2023)]{daunhawer2023identifiability}
I.~Daunhawer, A.~Bizeul, E.~Palumbo, A.~Marx, and J.~E. Vogt.
\newblock Identifiability results for multimodal contrastive learning.
\newblock \emph{ICLR}, 2023.

\bibitem[Zimmermann et al.(2021)]{zimmermann2021contrastive}
R.~S. Zimmermann, Y.~Sharma, S.~Schneider, M.~Bethge, and W.~Brendel.
\newblock Contrastive learning inverts the data generating process.
\newblock \emph{ICML}, 2021.

\bibitem[Brendel et al.(2023)]{brendel2023objectcentric}
W.~Brendel et al.
\newblock Object-centric learning and identifiability.
\newblock 2023. \todo{fix exact reference; also cite Kori et al.\ 2024}

\bibitem[Eastwood et al.(2023)]{eastwood2023dcies}
C.~Eastwood, A.~L. Nicolicioiu, J.~von K\"ugelgen, A.~Keki\'c, F.~Tr\"auble,
A.~Dittadi, and B.~Sch\"olkopf.
\newblock DCI-ES: An extended disentanglement framework with connections to
identifiability.
\newblock \emph{ICLR}, 2023.

\bibitem[Razavi et al.(2019)]{razavi2019vqvae2}
A.~Razavi, A.~van den Oord, and O.~Vinyals.
\newblock Generating diverse high-fidelity images with VQ-VAE-2.
\newblock \emph{NeurIPS}, 2019.

\end{thebibliography}

\end{document}
